\subsection{Naive cotangent complexes: source review and separate editorial arguments}\label{naive-cotangent-complexes-source-review-and-separate-editorial-arguments}

Authority: Stacks Project authors, commit \texttt{a04446e57ec1fbc252a871afcec7752fb2807b14}, \texttt{algebra.tex:35116–35804}. Original source: https://github.com/stacks/stacks-project/blob/a04446e57ec1fbc252a871afcec7752fb2807b14/algebra.tex\#L35116 . All statements below retain the original rings, presentation ideals, tensor actions and degree conventions. These are supporting editorial arguments, not replacement translation paragraphs or claims of novelty. Proposed local source fixes are recorded separately. A corrected translation must retain the source statement and its place in the original exposition.

Reading includes the entire section and reports OCC-00732 through OCC-00770. Bounded dependencies are \texttt{cotangent.tex:80–178,1830–2017}, \texttt{homology.tex:3605–3722}, \texttt{more-algebra.tex:8439–8555} and \texttt{simplicial.tex:6728–6775,6813–6849,6875–6991}. Additional read material in Cotangent 2019--2090 and More on Algebra 8556--8595 is lookahead, not an adjudicated chapter. A requested Simplicial read beyond its last line raised an index error after displaying the complete relevant examples; actual coverage, rather than the requested endpoint, is recorded. Two disk-index hits for ``cotangent complex Tor base change'' remain unread routing records. No PDF was used.

\subsubsection{Presentations and the exact homotopy maps}\label{presentations-and-the-exact-homotopy-maps}

For a ring map \(R\to S\), put \(P=R[S]\), with variable \([s]\) for every element of \(S\), and augmentation \(\epsilon([s])=s\). Its monomial basis includes the empty monomial. The ideal generated by \[
[s+t]-[s]-[t],\qquad [s][t]-[st],\qquad [r]-r
\] is exactly \(I=\ker\epsilon\): modulo these relations the map \(s\mapsto[s]\) is a unital ring map inverse to the induced augmentation. In the last expression \([r]\) denotes the variable of the image of \(r\) in \(S\). This also handles a noninjective base map and the zero ring.

For any polynomial presentation \(\alpha:P\twoheadrightarrow S\), with kernel \(I\), write \[
N(\alpha)_1=I/I^2,\quad N(\alpha)_0=\Omega_{P/R}\otimes_P S,
\qquad d(\overline i)=di\otimes1.
\] The product rule makes this well defined on \(I/I^2\). A map from its cokernel to an \(S\)-module \(M\) is exactly an \(R\)-derivation \(P\to M\) vanishing on \(I\), hence exactly an \(R\)-derivation \(S\to M\). Thus its cokernel is the original \(\Omega_{S/R}\), with its actual universal map. The degree-zero term is the free \(S\)-module on the differentials of the polynomial variables.

For the source square \(R\to R'\), \(\phi:S\to S'\), and a presentation lift \(\varphi:P\to P'\), the two maps are \[
\varphi_1(\overline i)=\overline{\varphi(i)},\qquad
\varphi_0(dp\otimes s)=d\varphi(p)\otimes\phi(s).
\] They are \(S\)-linear using \(\phi\) on the target and commute with \(d\). Each polynomial variable may be lifted through the surjection \(P'\to S'\), proving existence of \(\varphi\). Composition of these two formulas is exactly the formula for the composite lift.

For two lifts \(\varphi,\varphi'\), their difference takes values in \(I'\). The induced map \(D:P\to I'/(I')^2\) is an \(R\)-derivation because \[
(\varphi-\varphi')(pq)=\varphi(p)(\varphi-\varphi')(q)
+(\varphi-\varphi')(p)\varphi'(q),
\] and the two actions coincide modulo \((I')^2\). That action factors through \(S\). Therefore \[
h(dp\otimes s)=\phi(s)\,\overline{\varphi(p)-\varphi'(p)}
\] is a well-defined \(S\)-linear map \(N(\alpha)_0\to N(\alpha')_1\). On \(\overline i\) and \(dp\otimes s\), respectively, the formulas give \[
hd=\varphi_1-\varphi'_1,\qquad dh=\varphi_0-\varphi'_0.
\] If both base and target maps are isomorphisms, lift the reverse square and apply these identities to both composites and the identity lifts. This proves a homotopy equivalence, without any projectivity assertion about the conormal module. The reference to part (3) at line 35384 also needs part (2), which supplies these homotopies.

Taking the identity presentation of a polynomial algebra gives \((0\to\Omega_{S/R})\). Taking the zero-variable presentation of a surjection gives \((I/I^2\to0)\). The comparisons just proved are homotopy equivalences, not literal identifications of the canonical presentation complexes. In particular, for a nonzero ring \(T\), the canonical complex for \(T\to T\) has degree-zero term \(\bigoplus_{t\in T}T\,d[t]\), which is nonzero, although the complex is contractible. This justifies clarifying ``zero'' to ``homotopy equivalent to zero'' at line 35586 while retaining the original argument.

\subsubsection{The original truncation really is isomorphic}\label{the-original-truncation-really-is-isomorphic}

The standard resolution uses \(P_0=R[S]\), \(P_1=R[P_0]\), \(P_2=R[P_1]\), with the original, unnormalized associated chain complex \(L_n=\Omega_{P_n/R}\otimes_{P_n}S\). We do not replace this with a normalized complex. On a variable \([q]\in P_1\), the faces are \[
d_0([q])=q,\qquad d_1([q])=[\epsilon(q)].
\] On \([F]\in P_2\), they are \(F,[d_0F],[d_1F]\). These are the maps in Simplicial, \texttt{example-polynomial-algebra-maps}, lines 6875--6906. Let \(Q=\operatorname{coker}(L_2\to L_1)\), and denote the image in \(Q\) of \(d[q]\otimes1\) by \(e_q\). The exact relations defining \(Q\), one for each polynomial \(F\in R[P_0]\), are \[
\sum_q \epsilon_1\!\left(\frac{\partial F}{\partial[q]}\right)e_q
-e_{d_0F}+e_{d_1F}=0.\tag{1}
\] Every sum is finite and includes its original integer polynomial coefficients. Here \(\epsilon_1([q])=\epsilon(q)\).

Define \(u:L_1\to I/I^2\) by \[
u(d[q]\otimes1)=\overline{q-[\epsilon(q)]}.
\] More generally \(u(dF\otimes s)=s\overline{d_0F-d_1F}\). This is well defined because \(d_0-d_1\) is a derivation modulo \(I^2\), for the common action through \(S\), by the same product calculation as above. The simplicial identities give \((d_0-d_1)(d_0-d_1+d_2)=0\); explicitly the terms pair as \(d_0d_0=d_0d_1\), \(d_1d_0=d_0d_2\), and \(d_1d_1=d_1d_2\). Hence \(u\) kills all the relations (1) and descends to \(Q\). Its differential is exactly the degree-one differential into \(L_0\).

Here is an inverse, with no derived-category inference. First (1) with \(F=[[s]]\) gives \(e_{[s]}=0\). If \(i\in I\) and \(q\in P_0\), use \(F=[q+i]-[q]\). Its two faces are \(i\) and \(0\), so \[
e_{q+i}-e_q=e_i-e_0.\tag{2}
\] For \(i\in I,q\in P_0\), use \(F=[q][i]-[q][0]\). Its two faces are \(qi\) and \(0\), and the coefficients of \(e_q\) cancel. Thus \[
e_{qi}-e_0=\epsilon(q)(e_i-e_0).\tag{3}
\] Set \(v(i)=e_i-e_0\). Equation (2) makes \(v\) additive, and (3) makes it \(P_0\)-linear for the action through \(S\). It annihilates \(I^2\), and hence induces \(v:I/I^2\to Q\). Directly, \[
u(v(i))=\overline{i-[0]}-\overline{0-[0]}=\overline i.
\] Applying (2) with \(q=[\epsilon(q_0)]\), \(i=q_0-[\epsilon(q_0)]\), gives \[
v(u(e_{q_0}))=e_{q_0}-e_{[\epsilon(q_0)]}=e_{q_0}.
\] Thus \(u\) and \(v\) are inverse isomorphisms. Together with the identity on \(L_0\), they give the actual canonical isomorphism \[
\tau_{\leq1}L_{S/R}\ \xrightarrow{\ \cong\ }\ N(R[S]\to S).
\] This is the chain truncation of Homology 3694--3722: its degree-one term is precisely \(Q\). Cotangent 129--133 uses cochain degree \(-n\) for this same term, so its \(\tau_{\geq-1}\) is the same construction with that stated indexing. The weaker wording ``quasi-isomorphic'' in Cotangent 1937--1975 does not refute the stronger statement in Algebra 35170. No correction to that stronger statement is warranted. The explicit map agrees with Cotangent 1977--2017, and the proof here establishes its low-degree isomorphism directly, without asserting that all the sheaf-derived arguments in that chapter have been checked.

There is a different qualification needed at Algebra 35161. The standard augmentation is a homotopy equivalence of underlying simplicial sets, as Simplicial 6908--6991 explicitly states; it need not be a homotopy equivalence of simplicial \(R\)-algebras. For \(R=\mathbf Z\), \(S=\mathbf F_p\), a reverse simplicial algebra map would in degree zero give a unital map \(\mathbf F_p\to\mathbf Z[\mathbf F_p]\). This is impossible because \(p\cdot1\ne0\) in the polynomial ring. The actual set section sends \(s\) to the iterated bracket \([\ldots[s]\ldots]\). The extra degeneracy \(t_n(q)=[q]\) satisfies \(d_0t_n=\mathrm{id}\), \(d_{i+1}t_n=t_{n-1}d_i\), and \(s_{i+1}t_n=t_{n+1}s_i\), with the augmentation in degree \(-1\); these follow by deleting or inserting the indicated bracket. They contract the augmented underlying simplicial sets as in the cited construction. Adding ``on underlying simplicial sets'' retains the precise true claim.

\subsubsection{Jacobi-Zariski and the tensor comparison sequence}\label{jacobi-zariski-and-the-tensor-comparison-sequence}

Keep the source presentations \(P=A[x_s]\twoheadrightarrow B\), \(Q=B[y_t]\twoheadrightarrow C\), and \(T=A[x_s,y_t]\twoheadrightarrow C\), with kernels \(I,J,K\). Write \(W=IT\). Then \(T/W=Q\), \(J=K/W\), and \[
J/J^2=K/(K^2+W),\qquad
(I/I^2)\otimes_B C\cong W/KW.
\] For the latter formula, use \(T=P[y_t]\), so \(W=I\otimes_P T\), and then tensor over \(T\) with \(C=T/K\). Since \(I^2T\subset KW\), this gives the stated quotient. Thus the bottom row in the source is right exact, with map \(W/KW\to K/K^2\) induced by inclusion. No injectivity of that map is asserted. The top row is the split sequence of free \(C\)-modules on \(dx_s\), on \(dx_s,dy_t\), and on \(dy_t\), with their indicated polynomial tensor bases. Differentiation commutes with these maps.

For a cycle \(\overline j\in J/J^2\), choose a lift \(\overline k\in K/K^2\). The \(dy_t\) part of \(dk\) is zero, so \(dk\) comes from a unique vector \(v\in\bigoplus_s Cdx_s\). Send the cycle to the class of \(v\) in \(\Omega_{B/A}\otimes_B C\). Changing the lift by the image of \(W/KW\) changes \(v\) by the differential of that image's preimage; hence this boundary is well defined. A zero boundary allows subtraction of such a preimage from \(\overline k\), making \(\overline k\) a cycle. Conversely a lifted cycle has zero boundary. A cycle in \(K/K^2\) mapping to zero in \(J/J^2\) lifts from \(W/KW\), and injectivity of the top-left arrow forces that lift's differential to be zero. A vector from \(\bigoplus_s Cdx_s\) becomes a differential in the middle precisely when it arises from this boundary construction. Finally the right-exact cokernel sequence follows by lifting a differential on the right to the middle. These statements prove every exactness position of the displayed six-term sequence. Comparing presentations with the already proved homotopies supplies the canonical homology groups in the source.

The tensor proviso has an exact refinement. For any two-term complex \(N_1\xrightarrow d F\), where \(F\) is a free \(B\)-module, set \(H=\ker d\), \(M=\operatorname{im}d\), and \(\Omega=\operatorname{coker}d\). Tensor the exact sequences \(0\to M\to F\to\Omega\to0\) and \(0\to H\to N_1\to M\to0\) with a \(B\)-module \(C\). The first yields \[
\operatorname{Tor}_1^B(M,C)\cong\operatorname{Tor}_2^B(\Omega,C),
\quad
0\to\operatorname{Tor}_1^B(\Omega,C)\to M\otimes_B C\to F\otimes_B C.
\] Let \(\delta\) be the tensor connecting map from \(\operatorname{Tor}_1(M,C)\) to \(H\otimes C\). The surjection \(N_1\otimes C\to M\otimes C\) maps \(\ker(d\otimes1)\) onto \(\ker(M\otimes C\to F\otimes C)\). Its kernel is the image of \(H\otimes C\); the kernel of the map out of \(H\otimes C\) is \(\operatorname{im}\delta\). Consequently the natural exact sequence is \[
\operatorname{Tor}_2^B(\Omega,C)\xrightarrow\delta H\otimes_B C
\longrightarrow H_1(N\otimes_B C)
\longrightarrow\operatorname{Tor}_1^B(\Omega,C)\longrightarrow0.\tag{4}
\] In particular the comparison is an isomorphism exactly when \(\operatorname{Tor}_1(\Omega,C)=0\) and the particular map \(\delta\) is zero. Vanishing of the whole \(\operatorname{Tor}_2\) module, as assumed in the source, is sufficient. Formula (4) supplies the exact obstruction rather than silently weakening the source theorem's hypotheses in a translation.

For the source nullhomotopy, let \(\phi:B\to C\), \(P=A[B]\), \(Q=B[C]\), \(J=\ker(Q\to C)\). Compare \(f:P\to Q\), \([b]\mapsto[\phi(b)]\), with \(e:P\to Q\), \([b]\mapsto b\) as a coefficient. Their difference modulo \(J^2\) is a derivation, giving \[
h(d[b]\otimes b')=\phi(b')\,\overline{[\phi(b)]-b}.
\] Here the scalar \(b'\) acts via \(B\to C\). On \(I=\ker(P\to B)\), \(e(I)=0\), so \(hd=f_1\). Relative to \(B\), \(de(p)=0\), so \(dh=f_0\). This proves the composite is nullhomotopic. For a surjection \(A\to C\), applying the six-term sequence and the zero-variable presentations gives exactly \(I/I^2\to J/J^2\to\Omega_{B/A}\otimes_B C\to0\), retaining the original ideals.

\subsubsection{Base change with its full obstruction module}\label{base-change-with-its-full-obstruction-module}

Let \(R\to S\), \(P\twoheadrightarrow S\) with kernel \(I\), and an arbitrary ring map \(R\to R'\) be given. Put \(P'=P\otimes_R R'\), \(S'=S\otimes_R R'\), and \(I'=\ker(P'\to S')\). Since \(P\) is free over \(R\), tensoring \(0\to I\to P\to S\to0\) gives \[
0\to T:=\operatorname{Tor}_1^R(S,R')
\longrightarrow I\otimes_R R'\longrightarrow I'\longrightarrow0.
\] The image of \(I^2\otimes_R R'\) in \(I'\) is exactly \((I')^2\): products of sums of images of tensors are sums of such products, and coefficients in \(R'\) can be put in either tensor factor. Taking quotients therefore gives \[
T\longrightarrow (I/I^2)\otimes_R R'\longrightarrow I'/(I')^2\longrightarrow0.
\] Define \(O\) to be the image of the first map, with its induced \(S'\)-module structure. It is the exact kernel of the second map, not necessarily all of \(T\). The universal derivation identifies the degree-zero terms: \[
(\Omega_{P/R}\otimes_P S)\otimes_R R'
\cong\Omega_{P'/R'}\otimes_{P'}S'.
\] For polynomial \(P\), both sides are the free \(S'\)-module on the original differentials, and the map is their identity. An element of \(O\) has zero differential, since its image under this degree-zero isomorphism is the differential of zero. Thus the original comparison fits in a short exact sequence of complexes \[
0\to (O\to0)\to N(\alpha)\otimes_R R'\to N(\alpha')\to0.
\] Its full nontrivial homology assertion is \[
0\to O\to H_1(N(\alpha)\otimes_R R')
\to H_1(N(\alpha'))\to0,
\qquad H_0(N(\alpha)\otimes_R R')\cong H_0(N(\alpha')).\tag{5}
\] Hence the comparison is a quasi-isomorphism exactly when \(O=0\). Since its degree-zero map is an isomorphism and its degree-one map is surjective with kernel \(O\), this is also exactly when it is an isomorphism of these chosen presentation complexes. It suffices that \(\operatorname{Tor}_1^R(S,R')=0\); flatness of \(R'\) over \(R\) is unnecessary for that implication. Comparison with the canonical presentations gives the source's homotopy equivalence in this more general Tor-vanishing case. Tensoring an existing homotopy preserves its equations, without flatness.

For a strict extension of the flat-base hypothesis take \(R=\mathbf Z\), \(S=\mathbf Z[x]\), \(R'=\mathbf F_p\): \(S\) is free over \(R\), so the Tor group is zero, whereas \(R'\) is not flat (multiplication by \(p\) on \(\mathbf Z\) ceases to be injective after tensoring). For failure of unrestricted base change take \(R=\mathbf Z\), \(S=R'=\mathbf F_p\), and the zero-variable presentation \(P=\mathbf Z\), \(I=(p)\). The original complex after tensoring is \((\mathbf F_p\to0)\), while \(P'=S'=\mathbf F_p\) has zero kernel and zero complex. Here \(O=\mathbf F_p\), so the degree-one homology is genuinely lost. These examples verify both the stronger sufficient hypothesis and the necessary obstruction; they do not replace the flat-base theorem in the source.

For a directed system \(R_\lambda\to S_\lambda\), every polynomial and every equality between finite polynomial expressions occurs at some common stage. Hence \(R[S]=\operatorname{colim}R_\lambda[S_\lambda]\), and an element mapping to zero in \(S\) is already zero at a later stage, giving \(I=\operatorname{colim}I_\lambda\). Every element of \(I^2\) is a finite sum of finite products, so it and every relation in the quotient are similarly witnessed at a stage. The same argument applies to the free modules on the elements of \(S_\lambda\). It gives both terms, their scalar actions, and their differentials in the source's asserted colimit identity. Exactness of filtered colimits of modules follows by the same finite-element test, and will be used for localization below.

\subsubsection{Localization and why the reported splitting is valid}\label{localization-and-why-the-reported-splitting-is-valid}

For \(U\subset A\) multiplicative, flat base change to \(U^{-1}A\) compares \(N_{U^{-1}A/A}\) with the canonical identity complex over \(U^{-1}A\); the latter is contractible by the zero-variable presentation. Every term of the former is already a \(U^{-1}A\)-module, so tensoring it with \(U^{-1}A\) is an isomorphism. This proves the source localization complex is contractible. If \(U^{-1}A\to B\), localizing a presentation \(P\to B\) gives the presentation \(U^{-1}P\to B\). Its conormal module is \(U^{-1}(I/I^2)=I/I^2\), as all these scalars already act invertibly through \(B\); the differential module agrees for the same reason, by the polynomial differential formula. This proves the base-localization comparison with the original actions.

For the principal localization retain exactly \(\alpha:P\to B\), \(I\), \(f\in P\) mapping to \(g\in B\), and \(\beta:P[x]\to B_g\), \(x\mapsto g^{-1}\). Its kernel is \(J=(I,fx-1)\), because the indicated quotient is \(B[x]/(gx-1)=B_g\). In \(B[x]\), the polynomial \(gx-1\) is a nonzerodivisor: if \((gx-1)\sum_{i=0}^n b_ix^i=0\), the constant coefficient gives \(b_0=0\), and successive coefficients give \(b_i=gb_{i-1}=0\). Thus its conormal module is free on its class over \(B_g\).

The map \(\pi:P[x]\to(P/I^2)_f\), \(x\mapsto1/f\), sends \(J\) into the square-zero ideal \((I/I^2)_g\), kills \(J^2\), and induces a \(B_g\)-linear retraction of \((I/I^2)_g\to J/J^2\). On a fraction \(\overline i/g^n\) its composite is exactly \(\overline i/g^n\); hence the retraction proves injectivity for all fractions, including torsion. The conormal and differential rows consequently split as modules, and \[
J/J^2=(I/I^2)_g\oplus B_g\overline{fx-1},\quad
\Omega_{P[x]/A}\otimes_{P[x]}B_g
=(\Omega_{P/A}\otimes_P B_g)\oplus B_gdx.
\] Keep the full derivative \[
d(fx-1)=g^{-1}df+g\,dx=g\,(dx+g^{-2}df).
\] The change from \(dx\) to \(e=dx+g^{-2}df\) has inverse \(dx=e-g^{-2}df\), so it gives an isomorphism of the original complex with \[
N(\alpha)\otimes_B B_g\ \oplus\ (B_g\xrightarrow{g}B_g).
\] No \(df\) or denominator is discarded. The last summand contracts by multiplication by \(g^{-1}\) from degree zero to degree one. All formulas remain meaningful if \(B_g\) is the zero ring.

Report OCC-00761 incorrectly interprets ``Such a short exact sequence'' as referring to every sequence of complexes. In context the quotient is the displayed disk \((B_g\xrightarrow g B_g)\), where \(g\) is a unit. Any short exact sequence with that quotient splits in complexes: choose a lift \(v\) of its degree-one generator (the degree-one quotient is free), and define the lift of the degree-zero generator to be \(g^{-1}dv\). Its projection is the required generator and its differential condition is \(dv=g(g^{-1}dv)\). Extending linearly gives a chain section. This proves the general claim for exactly the displayed type of quotient, independently of the explicit splitting. The report is rejected and the true source claim is retained.

For arbitrary multiplicative \(U\subset B\), the rings \(B_g\), \(g\in U\), form a filtered system under divisibility and have colimit \(U^{-1}B\). Divisibility is a directed preorder, sufficient for the filtered-colimit argument; no antisymmetry assumption is required. The canonical comparison maps commute at every stage. Applying the proven principal result and filtered exactness gives the source's quasi-isomorphism. A stronger homotopy assertion is proved in the next section and retained as editorial material.

\subsubsection{Two-term quasi-isomorphisms and stronger localization}\label{two-term-quasi-isomorphisms-and-stronger-localization}

Let \(f:C\to D\) be a quasi-isomorphism of complexes in degrees one and zero, and suppose \(D_0\) is projective. The sequence \[
0\to C_1\xrightarrow{a=(f_1,-d_C)}D_1\oplus C_0
\xrightarrow{b=(d_D,f_0)}D_0\to0\tag{6}
\] is exact. Indeed \(a\) is injective by injectivity of \(H_1(f)\), and \(b\) is surjective by surjectivity of \(H_0(f)\). If \(d_Dv+f_0u=0\), injectivity of \(H_0(f)\) gives \(u=d_Cw\); then \(v+f_1w\) is a cycle. Surjectivity of \(H_1(f)\) supplies \(z\in\ker d_C\) with \(f_1z=v+f_1w\). The element \(z-w\) maps under \(a\) to \((v,u)\), proving exactness in the middle.

Choose a section \((s,t):D_0\to D_1\oplus C_0\) of \(b\), so \(d_Ds+f_0t=1\). Put \(g_0=t\). For \(v\in D_1\), the pair \((v-sd_Dv,-td_Dv)\) is in \(\ker b\); its unique lift under \(a\) defines \(g_1v\). Thus \[
f_1g_1=1-sd_D,\qquad d_Cg_1=td_D,
\] so \(g\) is a chain map and \(fg-1\) has homotopy \(-s\). For \(u\in C_0\), the pair \((-sf_0u,u-tf_0u)\) is in \(\ker b\); its unique lift defines \(hu\). Then \[
f_1h=-sf_0,\qquad d_Ch=tf_0-1.
\] Applying \(a\) to \(g_1f_1-1\) and \(hd_C\) gives the same two components, using the chain identities, so its injectivity gives \(g_1f_1-1=hd_C\). Thus \(gf-1=dh+hd\), and \(f\) is a homotopy equivalence. This proof imposes no projectivity on \(C_1,D_1,C_0\).

Apply this to the canonical arbitrary-localization map of the preceding section: the target degree-zero term is free over \(U^{-1}B\) on all its canonical variables. The established quasi-isomorphism is therefore a homotopy equivalence, even for an arbitrary multiplicative set. This strengthens the conclusion at Algebra 35694--35695. No compatible choices of homotopies at the filtered stages are assumed or needed. Keep this strengthened statement separate from the translation. It also shows how the source's later use of localization up to homotopy in More on Algebra 8585--8592 can be justified without silently treating all quasi-isomorphisms as homotopy equivalences.

\subsubsection{Stable conormal isomorphisms with both products}\label{stable-conormal-isomorphisms-with-both-products}

Retain the source complexes \(A_1\to A_0\), \(B_1\to B_0\), maps \(\varphi:A\to B\), \(\psi:B\to A\), and homotopies \[
1-\psi_1\varphi_1=hd_A,\quad1-\psi_0\varphi_0=d_Ah,
\quad1-\varphi_1\psi_1=h'd_B,\quad1-\varphi_0\psi_0=d_Bh'.
\] The displayed matrices have the exact types \[
T=\begin{pmatrix}\psi_1&h\\-d_B&\varphi_0\end{pmatrix}:
B_1\oplus A_0\to A_1\oplus B_0,
\qquad
U=\begin{pmatrix}\varphi_1&-h'\\d_A&\psi_0\end{pmatrix}:
A_1\oplus B_0\to B_1\oplus A_0.
\] Multiplying all four entries gives \[
TU=\begin{pmatrix}1&E\\0&1\end{pmatrix},\quad E=-\psi_1h'+h\psi_0,
\qquad
UT=\begin{pmatrix}1&F\\0&1\end{pmatrix},\quad F=\varphi_1h-h'\varphi_0.
\tag{7}
\] For example the lower-left entries are \(-d_B\varphi_1+\varphi_0d_A=0\) and \(d_A\psi_1-\psi_0d_B=0\), by the chain identities. The diagonal entries are the four displayed homotopy identities. Both shears are invertible, with off-diagonal entries \(-E\) and \(-F\). Consequently \[
T^{-1}=U(TU)^{-1}=(UT)^{-1}U,\qquad
U^{-1}=(TU)^{-1}T=T(UT)^{-1}.
\] Indeed each first expression is a right inverse and each second expression a left inverse; associativity gives equality of the left and right inverses. This proves the original stable isomorphism for arbitrary modules, not only finite free modules.

At line 35738 the source's \(-h'\circ\psi_1\) is ill typed; the required entry is \(-\psi_1\circ h'\). Report OCC-00766 correctly identifies that the displayed product alone is insufficient. For a concrete boundary, on \(V=\bigoplus_{n\geq0}ke_n\), let \(Se_n=e_{n+1}\), \(Le_0=0\), \(Le_{n+1}=e_n\). Then \(LS=1\), but \(SL(e_0)=0\ne e_0\). Thus neither factor need be invertible merely because one product is. The reverse calculation in (7) repairs the source proof without changing its theorem. The proposed local text only states that the reverse product too has identity diagonal; the complete computations remain here.

For two finite polynomial presentations of \(S/R\), the presentation homotopy equivalence and (7) give exactly \[
I/I^2\oplus S^{\oplus m}\cong J/J^2\oplus S^{\oplus n}.
\] For the source localization \(S_g\), use the proved homotopy equivalence from \(N(\alpha)\otimes_SS_g\) to the adjoined-variable presentation, and then to \(N(\beta)\). Applying the same two-term result gives exactly \[
(I/I^2)_g\oplus S_g^{\oplus m}\cong J/J^2\oplus S_g^{\oplus n}.
\] The reference required at line 35796 is \texttt{lemma-sum-two-terms}; \texttt{lemma-conormal-module} is the earlier statement for two polynomial presentations of the same algebra, whereas the tensor-localized complex in this argument need not itself arise from a polynomial presentation over \(R\). In particular, cancelling a free summand from a different stable isomorphism would be unjustified for arbitrary modules. The direct matrix comparison supplies the asserted result, including zero ranks and the zero ring.

\subsubsection{Propagation and bounded follow-up}\label{propagation-and-bounded-follow-up}

The presentation homotopy and the exact truncation comparison protect the source's genuine claims from proposed weakenings. The disk splitting rejects OCC-00761. The full tensor sequence (4), the obstruction (5), and the two-term inverse construction are separate consequences with zero translation operations. The original flat-base and quasi-isomorphism formulations remain identifiable. Previous conormal and base-change work in \texttt{ALGEBRA\_DIFFERENTIALS\_NOTES\_20260928.md} supplies the receiving comparison; the current arguments retain its actual ideals and module actions.

A source dependency also needs a separate follow-up: More on Algebra 8502 and 8504 write \((I/I^2)\otimes_A C\) where the exact quotient \(IT/IK\) requires \((I/I^2)\otimes_B C\). For \(A=k\), \(B=C=k[x]/(x^2)\), \(P=k[x]\), \(I=K=(x^2)\), and the zero-variable identity presentation \(B\to C\), the left expression with base \(k\) has dimension four over \(k\), whereas \(I/I^2\) and the required quotient have dimension two. The natural multiplication map is onto and has a two-dimensional kernel, so the printed identification is false. This explicit counterexample and the corrected map follow from the quotient calculation above. Record the two exact loci as a pending More on Algebra correction rather than editing an unadjudicated chapter or crediting its entire proof as checked. The cited Koszul-regular intersection theorem was not independently read in this batch.

The truncation isomorphism and the localization splitting need no corrective source operations. Copyediting preferences for hyphens, terse ``Follows from'' proofs, the causal ``being'', and two optional commas are not mathematical defects. The report about ``each \ldots{} are quasi-isomorphisms'' needs both verb and complement corrected to ``is a quasi-isomorphism''. Tensor-base clarifications retain the existing maps. Broader joint synthesis, literature comparison for any novelty claim, and a resulting short paper follow core completion. No chapter mutation, translation rewrite, TeX build, publication, or complete-project claim is made here.
